% GENERATED FILE. Edit canonical lessons, then run npm run generate:books.
% canonical-source-hash: fnv1a64:265389fa7984e1ef
% canonical-lessons: PA-R165-live-again, PA-R165-spend-again, PA-R165-save-again, PA-R165-sow-again

\chapter{Return to Living, Spending, Saving, and Sowing}
\label{ch:pa-four-familiar-actions-r4}
\hlchaptermodality{165}

\begin{chapteropening}
\textbf{By the end of this chapter:} \emph{I can retrieve four already-taught Punjabi action verbs after a long gap, one at a time, without claiming scored A1 writing or complete mock credit.}

Recall living, spending, saving, and sowing from short familiar cues without introducing new Punjabi.
\end{chapteropening}

\section[\texorpdfstring{\pa{ਜੀਉਣਾ} (jīuṇā)}{ਜੀਉਣਾ (jīuṇā)}]{An old word for living}

\label{lesson:PA-R165-live-again}



\begin{quote}
Hear \textbf{to live}. Recall \textbf{\pa{ਜੀਉਣਾ}} (\emph{jīuṇā}).
\end{quote}

\subsection*{Your turn}
Hear \textbf{\pa{ਜੀਉਣਾ}}. Recall \textbf{to live}. Then hear \textbf{to live} again. Say \textbf{\pa{ਜੀਉਣਾ}} before the model. Only this old word is needed.

\begin{tcolorbox}[breakable,colback=teal!4,colframe=teal!35!black,title={Before you move on}]
One last cue: \textbf{to live}. Recall \textbf{\pa{ਜੀਉਣਾ}}. Listen and retry once if it does not come back.
\end{tcolorbox}
\section[\texorpdfstring{\pa{ਖ਼ਰਚਣਾ} (kharcṇā)}{ਖ਼ਰਚਣਾ (kharcṇā)}]{An old word for spending}

\label{lesson:PA-R165-spend-again}



\begin{quote}
Hear \textbf{to spend}. Say \textbf{\pa{ਖ਼ਰਚਣਾ}} (\emph{kharcṇā}) before the model.
\end{quote}

\subsection*{Your turn}
Hear \textbf{\pa{ਖ਼ਰਚਣਾ}}. Recall \textbf{to spend}. Now hear \textbf{to spend}. Recall \textbf{\pa{ਖ਼ਰਚਣਾ}}. One old word is enough.

\begin{tcolorbox}[breakable,colback=teal!4,colframe=teal!35!black,title={Before you move on}]
Hear \textbf{to spend} once more. Say \textbf{\pa{ਖ਼ਰਚਣਾ}}. If needed, listen to the model and try again once.
\end{tcolorbox}
\section[\texorpdfstring{\pa{ਬਚਾਉਣਾ} (bacāuṇā)}{ਬਚਾਉਣਾ (bacāuṇā)}]{An old word for saving}

\label{lesson:PA-R165-save-again}



\begin{quote}
Hear \textbf{to save}. Recall \textbf{\pa{ਬਚਾਉਣਾ}} (\emph{bacāuṇā}).
\end{quote}

\subsection*{Your turn}
Hear \textbf{\pa{ਬਚਾਉਣਾ}}. Recall \textbf{to save}. Then hear \textbf{to save} again. Say \textbf{\pa{ਬਚਾਉਣਾ}} before the model. Keep to this one old meaning.

\begin{tcolorbox}[breakable,colback=teal!4,colframe=teal!35!black,title={Before you move on}]
One last cue: \textbf{to save}. Recall \textbf{\pa{ਬਚਾਉਣਾ}}. Listen and retry once if needed.
\end{tcolorbox}
\section[\texorpdfstring{\pa{ਬੀਜਣਾ} (bījṇā)}{ਬੀਜਣਾ (bījṇā)}]{An old word for sowing}

\label{lesson:PA-R165-sow-again}



\begin{quote}
Hear \textbf{to sow}. Recall \textbf{\pa{ਬੀਜਣਾ}} (\emph{bījṇā}).
\end{quote}

\subsection*{Your turn}
Pause after each familiar meaning, then hear its model:

\begin{itemize}
\raggedright
  \item \textbf{to save} — \textbf{\pa{ਬਚਾਉਣਾ}} (\emph{bacāuṇā}).
  \item \textbf{to live} — \textbf{\pa{ਜੀਉਣਾ}} (\emph{jīuṇā}).
  \item \textbf{to sow} — \textbf{\pa{ਬੀਜਣਾ}} (\emph{bījṇā}).
  \item \textbf{to spend} — \textbf{\pa{ਖ਼ਰਚਣਾ}} (\emph{kharcṇā}).
\end{itemize}

Only the order of these old words has changed.

\begin{tcolorbox}[breakable,colback=teal!4,colframe=teal!35!black,title={Before you move on}]
Hear \textbf{to sow}, \textbf{to spend}, \textbf{to live}, then \textbf{to save}. Pause after each cue and recall its Punjabi word before the model: \textbf{\pa{ਬੀਜਣਾ}}, \textbf{\pa{ਖ਼ਰਚਣਾ}}, \textbf{\pa{ਜੀਉਣਾ}}, \textbf{\pa{ਬਚਾਉਣਾ}}. Revisit only the old word you missed.
\end{tcolorbox}
